%\vspace{1.5\baselineskip}
\begin{figure}[H]
    \begin{center}
        \begin{circuitikz}[european]
            \ctikzset{tripoles/mos style/arrows}[line width=thick, european]
            % Widerstand
            \draw (-3,5) coordinate(resistor-begin) to[resistor, l={$R_0$}] (-3,0.5) coordinate(resistor-end);

            \draw (resistor-begin) node[circle, fill, inner sep=1.5pt]{};
            
            % NMOS Transistor
            \draw (-3,-0.5) node[nmos, anchor=drain] (nmosL) {$N_0$};

            % Verbindung von Drain des NMOS mit Widerstand
            \draw (resistor-end) -- (nmosL.drain);
            \draw (nmosL.gate) -- ++(-0.25,0) |- (resistor-end)
                node[circle, fill, inner sep=1.5pt]{}
                node[above, midway]{$V_{ref}$};
            %
            %*************************************************************************
            %
            % Widerstand R
            \draw (0,5) coordinate (resistor-begin1) to[resistor, l={$R_1$}] (0,2) coordinate (resistor-end1);
            \draw (0,2) coordinate (resistor-begin2) to[resistor, l={$R_2$}] (0,-1) coordinate (resistor-end2);
            \draw (0,-1) coordinate (resistor-begin3) to[resistor, l={$R_3$}] (0,-4) coordinate (resistor-end3);

            % NMOSR 
            \draw (3,2) node[nmos, rotate=270] (nmosR1) {}
                node at (nmosR1.center) [anchor=center, yshift=-10pt] {$N_1$}
                (nmosR1.source) -- (resistor-end1) % Verbinung zu Wiederständen 
                node[circle, fill, inner sep=1.5pt]{}; % Punkt am Ende der Verbindung 
            \draw (3,-1) node[nmos, rotate=270] (nmosR2) {}
                node at (nmosR2.center) [anchor=center, yshift=-10pt] {$N_2$}
                (nmosR2.source) -- (resistor-end2)
                node[circle, fill, inner sep=1.5pt]{};

            % NMOS Drain Verbindung
            \draw (nmosR1.drain) -- ++ (0.25,0) % horizontal nach rechts
                -- ++ (0,-3) coordinate (comp) % vertikale Linie 
                node[circle, fill, inner sep=1.5pt]{};

            \draw (nmosR1.gate) -- ++(0,0.25) -- ++(-0.75,0) 
                node[above, midway]{$V_{res}$};
            \draw (nmosR2.gate) -- ++(0,0.25) -- ++(-0.75,0)
                node[above, midway]{$V_{res0}$};

            % Dreieck am Ende von Vcomp
            \draw (6.75,-0.51) node[op amp] (triangle){}
                (triangle.+) -- (nmosR2.drain);
            \draw (triangle) node[above, yshift=27pt] {Komparator};  

            \draw (triangle.-) -- ++(-0.5,0) node[above]{$V_{ref}$};

            % Verbindung Ground
            \draw (nmosL.source) |- (resistor-end3)
                node[circle, fill, inner sep=1.5pt]{}
                -- ++(0,-0.25) node[sground]{}
                node[right, xshift=5pt, yshift=-5pt]{Gnd};     

            % Verbindung VDD
            \draw (resistor-begin1) -- ++(-3.73,0)
                node[above] {$V_{DD}$};

            % Arrow 
            \draw[->, thick, shorten >=10pt, shorten <=10pt, >=stealth] (resistor-end1) to[out=225,in=135] (resistor-end3)
                node[midway, left, xshift=-35pt, yshift=-10pt]{$V_{DD1}$};
                
            \draw[->, thick, shorten >=10pt, shorten <=10pt, >=stealth] (resistor-end2) to[out=-25,in=25] (resistor-end3)
                node[right,xshift=20pt, yshift=50pt]{$V_{DD2}$};
        \end{circuitikz}
    \end{center}
    \caption[Generierung des Referenzstrom durch Widerstand]{\\Quelle: Generierung des Referenzstrom durch Widerstand (Eigene Darstellung)}
    \label{fig:Vref_old}
\end{figure}
%\vspace{1.5\baselineskip}